\makeatletter
\def\input@path{{./}{../}{../../}{../../../}{../../../../}{preamble/}{../preamble/}{../../preamble/}{../../../preamble/}}
\makeatother

\input{preamble_exams.tex}

\title{Decision-Making Under Uncertainty: Exam Solutions}
\author{}
\date{}

\begin{document}
	\maketitle
	\vspace{-2cm}
	\setcounter{tocdepth}{1}
	\setcounter{secnumdepth}{2}
	\phantomsection
	\label{toc}
	\tableofcontents

	\vspace{-5mm}
	\section*{Learning objective}\vspace{-5mm}
	Apply minimax regret and Expected Monetary Value (EMV) to select an alternative under uncertainty, and use sensitivity analysis to determine how a subjective probability changes the preferred alternative.

	\section*{Criteria assessment}\vspace{-5mm}
	This task sheet assesses Criteria C5 through C7. Each problem assesses the criterion indicated in its title, and each problem is worth 2 points.
	Partial credit may be awarded when the correct procedure is clearly shown but arithmetic errors are made. Answers that do not include the required procedure will not receive credit.
	Failure to write your name on the task sheet or submitting it after the deadline will result in a one-point deduction from each criterion.

	A minimum of 4 points is required to pass each criterion.

	\section*{Evaluated criteria}
	\begin{enumerate}[label=\textbf{C\arabic*}, start=5]
		\item Applies the Minimax Regret criterion by constructing an opportunity-loss table and selecting the alternative with the smallest maximum regret.
		\item Calculates EMV using given probabilities and selects the alternative with the highest expected payoff.
		\item Uses subjective probabilities and sensitivity analysis to identify threshold probabilities and the preferred alternative over each feasible range.
	\end{enumerate}

	\sectionScored{Artisan Apiary Expansion}{C5-\sml{2}}{sec:garden-c5}{
		An artisan apiary must choose one of three hive-expansion plans. Annual net benefits (in thousands of dollars) depend on the honey season.
		\begin{center}\small\setlength{\tabcolsep}{3pt}\renewcommand{\arraystretch}{1.15}
			\begin{tabular}{lccc}\toprule
			Alternative & Excellent & Average & Poor\\ \midrule
			Medium&58&49&23\\
			Large&76&34&-14\\
			Small&38&43&35\\ \bottomrule
			\end{tabular}
		\end{center}\vspace{-3mm}
		Use the Minimax Regret criterion to determine which plan should be selected. Show the best payoff in each state, the regret table, and each alternative's maximum regret. \score{2}

		\subsection*{C5 --- Solution}
		The best payoff in each state is
		\[
		b(S_1)=\max\{76,58,38\}=76,\quad
		b(S_2)=\max\{34,49,43\}=49,\quad
		b(S_3)=\max\{-14,23,35\}=35.
		\]
		Using $r(A_i,S_j)=b(S_j)-P_{ay}(A_i,S_j)$ gives
		\begin{center}\renewcommand{\arraystretch}{1.25}\begin{tabular}{lcccc}\toprule
				& $S_1$ & $S_2$ & $S_3$ & Maximum regret\\ \midrule
				$A_1$&$76-58=18$&$49-49=0$&$35-23=12$&$18$\\
				$A_2$&$76-76=0$&$49-34=15$&$35-(-14)=49$&$49$\\
				$A_3$&$76-38=38$&$49-43=6$&$35-35=0$&$38$\\ \bottomrule\end{tabular}\end{center}
		Thus $\min\{18,49,38\}=18$, so the Minimax Regret criterion selects $\boxed{A_1}$, the medium expansion.
		\secInicio[sec:c5-exam-solutions][sec:garden-c5]
	}

	\sectionScored{Regional Theater Capacity}{C5-\sml{2}}{sec:ferry-c5}{
		A theater operator is choosing a seating-capacity plan. Seasonal profit (in thousands of dollars) is shown below; no reliable probabilities are available.
		\begin{center}\setlength{\tabcolsep}{2pt}\renewcommand{\arraystretch}{1.15}
			\begin{tabular}{lccc}\toprule
			Alternative & High & Normal & Low\\ \midrule
			Add two seating sections&108&40&-28\\
			Lease a nearby hall&59&52&34\\
			Add one seating section&82&61&14\\
			Current seating&34&34&34\\ \bottomrule
			\end{tabular}
		\end{center}\vspace{-3mm}
		Determine the decision using Minimax Regret. \score{2}

		\subsection*{C5 --- Solution}
		The statewise best payoffs are $b(S_1)=108$, $b(S_2)=61$, and $b(S_3)=34$.
		\begin{center}\renewcommand{\arraystretch}{1.25}\begin{tabular}{lcccc}\toprule
				& $S_1$ & $S_2$ & $S_3$ & Maximum regret\\ \midrule
				$A_1$&$108-108=0$&$61-40=21$&$34-(-28)=62$&$62$\\
				$A_2$&$108-59=49$&$61-52=9$&$34-34=0$&$49$\\
				$A_3$&$108-82=26$&$61-61=0$&$34-14=20$&$26$\\
				$A_4$&$108-34=74$&$61-34=27$&$34-34=0$&$74$\\ \bottomrule\end{tabular}\end{center}
		Since $\min\{62,49,26,74\}=26$, the recommended decision is $\boxed{A_3}$: add one seating section.
		\secInicio[sec:c5-exam-solutions][sec:ferry-c5]
	}

	\sectionScored{Emergency Shelter Staffing}{C5-\sml{2}}{sec:clinic-c5}{
		A relief organization must select an emergency-shelter staffing plan. Net community value (in thousands of dollars) depends on requests for shelter.
		\begin{center}\setlength{\tabcolsep}{2pt}\renewcommand{\arraystretch}{1.15}
			\begin{tabular}{lcccc}\toprule
			Alternative & Very high & High & Moderate & Low\\ \midrule
			Flexible&74&64&49&12\\
			Full&96&54&24&-24\\
			Limited&48&51&46&38\\ \bottomrule
			\end{tabular}
		\end{center}\vspace{-3mm}
		Use Minimax Regret to select a staffing plan. \score{2}

		\subsection*{C5 --- Solution}
		The best-payoff vector is
		\[
		(b(S_1),b(S_2),b(S_3),b(S_4))=(96,64,49,38).
		\]
		\begin{center}\renewcommand{\arraystretch}{1.25}\begin{tabular}{lccccc}\toprule
				&$S_1$&$S_2$&$S_3$&$S_4$&Maximum regret\\ \midrule
				$A_1$&22&0&0&26&26\\
				$A_2$&0&10&25&62&62\\
				$A_3$&48&13&3&0&48\\ \bottomrule\end{tabular}\end{center}
		For example, $r(A_2,S_4)=38-(-24)=62$ and $r(A_3,S_2)=64-51=13$. The smallest maximum regret is $\min\{26,62,48\}=26$, so select $\boxed{A_1}$, the flexible staffing plan.
		\secInicio[sec:c5-exam-solutions][sec:clinic-c5]
	}

	\sectionScored{Coffee Roastery Product Mix}{C6-\sml{2}}{sec:farm-c6}{
		A coffee roastery is selecting a product mix. Profits are in thousands of dollars.
		\begin{center}\setlength{\tabcolsep}{2pt}\renewcommand{\arraystretch}{1.15}
			\begin{tabular}{lccc}\toprule
			&Strong&Stable&Weak\\
			$P(S_j)$&0.45&0.35&0.20\\ \midrule
			Single-origin&86&44&-12\\
			Balanced&64&58&22\\
			Classic&42&47&37\\ \bottomrule
			\end{tabular}
		\end{center}\vspace{-3mm}
		Calculate every EMV and apply the Bayes decision criterion. \score{2}

		\subsection*{C6 --- Solution}
		First, $0.45+0.35+0.20=1$. Then
		\begin{align*}
			EMV(A_1)&=(0.45)(86)+(0.35)(44)+(0.20)(-12)=38.7+15.4-2.4=51.7,\\
			EMV(A_2)&=(0.45)(64)+(0.35)(58)+(0.20)(22)=28.8+20.3+4.4=53.5,\\
			EMV(A_3)&=(0.45)(42)+(0.35)(47)+(0.20)(37)=18.9+16.45+7.4=42.75.
		\end{align*}
		The highest EMV is $\max\{51.7,53.5,42.75\}=53.5$. Therefore, choose $\boxed{A_2}$, the balanced product mix, with expected profit $\$53{,}500$.
		\secInicio[sec:c6-exam-solutions][sec:farm-c6]
	}

	\sectionScored{Aquarium Program Scale}{C6-\sml{2}}{sec:museum-c6}{
		An aquarium is deciding the scale of a temporary education program. Payoffs are in thousands of dollars.
		\begin{center}\setlength{\tabcolsep}{2pt}\renewcommand{\arraystretch}{1.15}
			\begin{tabular}{lccc}\toprule
			&High&Expected&Low\\
			$P(S_j)$&0.30&0.40&0.30\\ \midrule
			Major&98&38&-24\\
			Compact&52&47&34\\
			Standard&73&53&18\\
			Regular tours only&27&27&27\\ \bottomrule
			\end{tabular}
		\end{center}\vspace{-3mm}
		Calculate the EMVs and recommend an alternative. \score{2}

		\subsection*{C6 --- Solution}
		The probabilities are valid because $0.30+0.40+0.30=1$.
		\begin{align*}
			EMV(A_1)&=(0.30)(98)+(0.40)(38)+(0.30)(-24)=29.4+15.2-7.2=37.40,\\
			EMV(A_2)&=(0.30)(52)+(0.40)(47)+(0.30)(34)=15.6+18.8+10.2=44.60,\\
			EMV(A_3)&=(0.30)(73)+(0.40)(53)+(0.30)(18)=21.9+21.2+5.4=48.50,\\
			EMV(A_4)&=(0.30)(27)+(0.40)(27)+(0.30)(27)=27.00.
		\end{align*}
		Because $48.50$ is the highest EMV, the Bayes decision is $\boxed{A_3}$, the standard program, with expected payoff $\$48{,}500$.
		\secInicio[sec:c6-exam-solutions][sec:museum-c6]
	}

	\sectionScored{Wind-Farm Maintenance Staffing}{C6-\sml{2}}{sec:solar-c6}{
		A wind-farm operator is choosing a maintenance staffing plan for next quarter. Payoffs are in thousands of dollars.
		\begin{center}\setlength{\tabcolsep}{1.5pt}\renewcommand{\arraystretch}{1.15}
			\begin{tabular}{lcccc}\toprule
			&Surge&Strong&Steady&Slow\\
			$P(S_j)$&0.20&0.30&0.35&0.15\\ \midrule
			Moderate hiring&79&68&48&12\\
			Rapid hiring&105&64&24&-32\\
			Current staff&53&55&51&42\\ \bottomrule
			\end{tabular}
		\end{center}\vspace{-3mm}
		Use EMV to determine the preferred staffing plan. \score{2}

		\subsection*{C6 --- Solution}
		The probability check is $0.20+0.30+0.35+0.15=1$.
		\begin{align*}
			EMV(A_1)&=(0.20)(79)+(0.30)(68)+(0.35)(48)+(0.15)(12)=15.8+20.4+16.8+1.8=54.8,\\
			EMV(A_2)&=(0.20)(105)+(0.30)(64)+(0.35)(24)+(0.15)(-32)=21+19.2+8.4-4.8=43.8,\\
			EMV(A_3)&=(0.20)(53)+(0.30)(55)+(0.35)(51)+(0.15)(42)=10.6+16.5+17.85+6.3=51.25.
		\end{align*}
		Thus $\max\{54.8,43.8,51.25\}=54.8$, so choose $\boxed{A_1}$, moderate hiring, with expected payoff $\$54{,}800$.
		\secInicio[sec:c6-exam-solutions][sec:solar-c6]
	}

	\sectionScored{Concert Merchandise Order}{C7-\sml{2}}{sec:catering-c7}{
		A concert vendor must choose between a large merchandise order ($A_1$) and a small merchandise order ($A_2$). Payoffs (in thousands of dollars) depend on whether concert attendance is high ($S_1$) or low ($S_2$).
		\begin{center}\setlength{\tabcolsep}{4pt}\renewcommand{\arraystretch}{1.15}
			\begin{tabular}{lcc}\toprule
			Alternative&High ($S_1$)&Low ($S_2$)\\ \midrule
			Large ($A_1$)&80&20\\
			Small ($A_2$)&50&40\\ \bottomrule
			\end{tabular}
		\end{center}\vspace{-3mm}
		Let $P(S_1)=p$ and $P(S_2)=1-p$. Find the threshold probability and state the preferred alternative for every $0\le p\le1$. \score{2}

		\subsection*{C7 --- Solution}
		\begin{align*}
			EMV(A_1)&=80p+20(1-p)=20+60p,\\
			EMV(A_2)&=50p+40(1-p)=40+10p.
		\end{align*}
		At the threshold, $20+60p=40+10p$, so $50p=20$ and $p=\frac{2}{5}=0.4$. Since $A_2$ has the higher EMV at $p=0$ and $A_1$ has the higher EMV at $p=1$,
		\[
		\boxed{\begin{cases}
				A_2,&0\le p<\frac{2}{5},\\
				A_1\text{ and }A_2\text{ are tied},&p=\frac{2}{5},\\
				A_1,&\frac{2}{5}<p\le1.
		\end{cases}}
		\]
		\secInicio[sec:c7-exam-solutions][sec:catering-c7]
	}

	\sectionScored{Meal-Kit Launch}{C7-\sml{2}}{sec:bike-c7}{
		A meal-kit company is comparing a bold launch ($A_1$) with a measured launch ($A_2$). Payoffs are in thousands of dollars.
		\begin{center}\setlength{\tabcolsep}{2pt}\renewcommand{\arraystretch}{1.15}
			\begin{tabular}{lccc}\toprule
				Alternative&High&Moderate&Low\\ \midrule
				Bold ($A_1$)&60&30&0\\
				Measured ($A_2$)&35&45&30\\ \bottomrule
			\end{tabular}
		\end{center}\vspace{-3mm}
		Suppose $P(S_1)=p$ and $P(S_2)=\frac{p}{3}$. Determine the feasible range of $p$, the threshold, and the preferred launch over the entire range. \score{2}

		\subsection*{C7 --- Solution}
		Because the three state probabilities must sum to one,
		\[
		P(S_3)=1-p-\frac{p}{3}=1-\frac{4p}{3}.
		\]
		Nonnegative probabilities require $p\ge0$ and $1-\frac{4p}{3}\ge0$, so the feasible range is $0\le p\le\frac{3}{4}$.
		\begin{align*}
			EMV(A_1)&=60p+30\left(\frac{p}{3}\right)+0\left(1-\frac{4p}{3}\right)=70p,\\
			EMV(A_2)&=35p+45\left(\frac{p}{3}\right)+30\left(1-\frac{4p}{3}\right)=30+10p.
		\end{align*}
		Equating the EMVs gives $70p=30+10p$, hence
		\[
		60p=30\quad\Rightarrow\quad p=\frac{1}{2}.
		\]
		Therefore,
		\[
		\boxed{\begin{cases}
				A_2,&0\le p<\frac{1}{2},\\
				A_1\text{ and }A_2\text{ are tied},&p=\frac{1}{2},\\
				A_1,&\frac{1}{2}<p\le\frac{3}{4}.
		\end{cases}}
		\]
		\secInicio[sec:c7-exam-solutions][sec:bike-c7]
	}

	\sectionScored{Publishing-Platform Bid}{C7-\sml{2}}{sec:rainwater-c7}{
		A publishing firm must choose between a standard platform bid ($A_1$) and an experimental platform bid ($A_2$). Payoffs are in thousands of dollars.
		\begin{center}\setlength{\tabcolsep}{2pt}\renewcommand{\arraystretch}{1.15}
			\begin{tabular}{lccc}\toprule
				Alternative&Immediate&Revision&Delayed\\ \midrule
				Standard ($A_1$)&30&70&70\\
				Experimental ($A_2$)&50&50&30\\ \bottomrule
			\end{tabular}
		\end{center}\vspace{-3mm}
		Suppose $P(S_1)=3p$ and $P(S_2)=p$. Find the feasible range, threshold probability, and preferred bid for each range of $p$. \score{2}

		\subsection*{C7 --- Solution}
		The remaining probability is
		\[
		P(S_3)=1-3p-p=1-4p.
		\]
		The probability conditions $p\ge0$ and $1-4p\ge0$ give the feasible range $0\le p\le\frac{1}{4}$.
		\begin{align*}
			EMV(A_1)&=30(3p)+70p+70(1-4p)=70-120p,\\
			EMV(A_2)&=50(3p)+50p+30(1-4p)=30+80p.
		\end{align*}
		At indifference,
		\[
		70-120p=30+80p
		\quad\Rightarrow\quad 40=200p
		\quad\Rightarrow\quad p=\frac{1}{5}=0.2.
		\]
		Here $A_1$ is preferred below the threshold, while the faster-growing EMV of $A_2$ makes it preferred above the threshold:
		\[
		\boxed{\begin{cases}
				A_1,&0\le p<\frac{1}{5},\\
				A_1\text{ and }A_2\text{ are tied},&p=\frac{1}{5},\\
				A_2,&\frac{1}{5}<p\le\frac{1}{4}.
		\end{cases}}
		\]
		\secInicio[sec:c7-exam-solutions][sec:rainwater-c7]
	}


\section{Summarized Solutions}

\subsection{Artisan Apiary Expansion \hfill C5--2}
The best-payoff vector is $(76,49,35)$. The regret table and row maxima are
\[
\begin{array}{c|ccc|c}&S_1&S_2&S_3&\max r\\ \hline A_1&18&0&12&18\\ A_2&0&15&49&49\\ A_3&38&6&0&38\end{array}
\]
Thus $\min\{18,49,38\}=18$, so select $\boxed{A_1}$, the medium expansion.

\subsection{Regional Theater Capacity \hfill C5--2}
The best-payoff vector is $(108,61,34)$. The regret table and row maxima are
\[
\begin{array}{c|ccc|c}&S_1&S_2&S_3&\max r\\ \hline A_1&0&21&62&62\\ A_2&49&9&0&49\\ A_3&26&0&20&26\\ A_4&74&27&0&74\end{array}
\]
Thus $\min\{62,49,26,74\}=26$, so select $\boxed{A_3}$, add one seating section.

\subsection{Emergency Shelter Staffing \hfill C5--2}
The best-payoff vector is $(96,64,49,38)$. The regret table and row maxima are
\[
\begin{array}{c|cccc|c}&S_1&S_2&S_3&S_4&\max r\\ \hline A_1&22&0&0&26&26\\ A_2&0&10&25&62&62\\ A_3&48&13&3&0&48\end{array}
\]
Thus $\min\{26,62,48\}=26$, so select $\boxed{A_1}$, flexible staffing.

\subsection{Coffee Roastery Product Mix \hfill C6--2}
Since $0.45+0.35+0.20=1$,
\begin{align*}
EMV(A_1)&=(0.45)(86)+(0.35)(44)+(0.20)(-12)=51.7,\\
EMV(A_2)&=(0.45)(64)+(0.35)(58)+(0.20)(22)=53.5,\\
EMV(A_3)&=(0.45)(42)+(0.35)(47)+(0.20)(37)=42.75.
\end{align*}
The largest EMV is $53.5$, so select $\boxed{A_2}$, the balanced product mix ($\$53{,}500$).

\subsection{Aquarium Program Scale \hfill C6--2}
\begin{align*}EMV(A_1)&=37.40,&EMV(A_2)&=44.60,\\ EMV(A_3)&=48.50,&EMV(A_4)&=27.00.\end{align*}
Using probabilities $(0.30,0.40,0.30)$, the largest EMV is $48.50$, so select $\boxed{A_3}$, the standard program ($\$48{,}500$).

\subsection{Wind-Farm Maintenance Staffing \hfill C6--2}
\begin{align*}EMV(A_1)&=(0.20)(79)+(0.30)(68)+(0.35)(48)+(0.15)(12)=54.8,\\ EMV(A_2)&=(0.20)(105)+(0.30)(64)+(0.35)(24)+(0.15)(-32)=43.8,\\ EMV(A_3)&=(0.20)(53)+(0.30)(55)+(0.35)(51)+(0.15)(42)=51.25.\end{align*}
The largest EMV is $54.8$, so select $\boxed{A_1}$, moderate hiring ($\$54{,}800$).

\subsection{Concert Merchandise Order \hfill C7--2}
\[
EMV(A_1)=20+60p,\qquad EMV(A_2)=40+10p.
\]
Indifference requires $20+60p=40+10p$, so $p=\frac25=0.4$. Therefore
\[
\boxed{A_2\ (0\le p<\tfrac25);\quad A_1=A_2\ (p=\tfrac25);\quad A_1\ (\tfrac25<p\le1).}
\]

\subsection{Meal-Kit Launch \hfill C7--2}
Here $P(S_3)=1-\frac{4p}{3}$, so $0\le p\le\frac34$. Also,
\[
EMV(A_1)=70p,\qquad EMV(A_2)=30+10p.
\]
Indifference requires $70p=30+10p$, so $p=\frac12$. Therefore
\[
\boxed{A_2\ (0\le p<\tfrac12);\quad A_1=A_2\ (p=\tfrac12);\quad A_1\ (\tfrac12<p\le\tfrac34).}
\]

\subsection{Publishing-Platform Bid \hfill C7--2}
Here $P(S_3)=1-4p$, so $0\le p\le\frac14$. Also,
\[
EMV(A_1)=70-120p,\qquad EMV(A_2)=30+80p.
\]
Indifference requires $70-120p=30+80p$, so $p=\frac15=0.2$. Therefore
\[
\boxed{A_1\ (0\le p<\tfrac15);\quad A_1=A_2\ (p=\tfrac15);\quad A_2\ (\tfrac15<p\le\tfrac14).}
\]

\end{document}
